% Generated draft: requires mathematical and source review.
\documentclass[11pt]{article}
\usepackage[margin=25mm]{geometry}
\usepackage{fontspec}
\setmainfont{Latin Modern Roman}
\usepackage{amsmath,amssymb,mathtools,unicode-math}
\setmathfont{Latin Modern Math}
\usepackage{enumitem,needspace,xurl,hyperref}
\setlength{\parindent}{0pt}
\setlength{\parskip}{5pt}
\setlength{\emergencystretch}{3em}
\newcommand{\N}{\mathbb{N}}
\newcommand{\Z}{\mathbb{Z}}
\newcommand{\Q}{\mathbb{Q}}
\newcommand{\R}{\mathbb{R}}
\newcommand{\C}{\mathbb{C}}
\newcommand{\F}{\mathbb{F}}
\newcommand{\A}{\mathbb{A}}
\newcommand{\PP}{\mathbb{P}}
\newcommand{\E}{\mathbb{E}}
\newcommand{\eps}{\varepsilon}
\newcommand{\dd}{\,\mathrm{d}}
\begin{document}
\section*{UnsolvedMath ID: 700923}
\textbf{Source problem code:} catalogue entry
\subsection{Definitions and mathematical statement}
A group is periodic if every element has finite order. It is locally soluble if each of its finitely generated subgroups is soluble, meaning that its derived series reaches the trivial subgroup after finitely many steps.

Groups \(A\) and \(B\) are locally equivalent if every finitely generated subgroup of \(A\) is isomorphic to a subgroup of \(B\), and every finitely generated subgroup of \(B\) is isomorphic to a subgroup of \(A\). A group \(G\) is categorical, in the terminology of this problem, if
\[
\forall\text{ groups }H,\qquad H\text{ locally equivalent to }G\ \Longrightarrow\ H\cong G.
\]
The quantifier over \(H\) is unrestricted; the source does not impose periodicity or local solubility on \(H\).

For a prime \(p\), a \(p\)-subgroup is a subgroup whose elements all have order a power of \(p\), and a Sylow \(p\)-subgroup means a maximal \(p\)-subgroup under inclusion. A Chernikov group is a group \(C\) having a normal subgroup \(D\) of finite index with
\[
D\cong\prod_{i=1}^{r} C_{p_i^\infty},\qquad 0\leq r<\infty,
\]
where the \(p_i\) are primes, repetitions are permitted, and \(C_{p^\infty}\) denotes the quasicyclic \(p\)-group, the direct limit of cyclic groups of orders \(p,p^2,p^3,\ldots\) under their standard injective inclusions. The empty product is trivial, so finite groups are included.

The source calls a group hyperfinite if it has a well ordered ascending normal series with finite factors. Under the convention that every term of a normal series is normal in the whole group, this means that there are an ordinal \(\lambda\) and subgroups \((G_\alpha)_{\alpha\leq\lambda}\) such that
\[
G_0=1,\qquad G_\lambda=G,\qquad
G_\alpha\leq G_\beta\quad(\alpha\leq\beta),\qquad
G_\alpha\trianglelefteq G\quad(\alpha\leq\lambda),
\]
\[
G_\delta=\bigcup_{\alpha<\delta}G_\alpha
\quad\text{for every nonzero limit ordinal }\delta\leq\lambda,
\qquad
|G_{\alpha+1}/G_\alpha|<\infty\quad(\alpha<\lambda).
\]
This normality convention is an explicit interpretation: the supplied record does not say whether all terms must be normal in \(G\), or only each term in its successor. That distinction requires confirmation against the cited source.

The question attributed to V. V. Belyaev is whether, for every periodic locally soluble group \(G\),
\[
G\text{ is categorical}
\quad\Longleftrightarrow\quad
\left(\begin{gathered}
G\text{ is hyperfinite, and}\\
\text{for every prime }p,\text{ every Sylow }p\text{-subgroup of }G
\text{ is Chernikov}
\end{gathered}\right).
\]
The interpretation of hyperfinite in this display is subject to the normal-series qualification above.
\subsection{Short English statement}
Belyaev asks whether a periodic locally soluble group is determined up to isomorphism by the isomorphism types of its finitely generated subgroups exactly when it is hyperfinite and all its Sylow subgroups are Chernikov. The supplied normal-series terminology requires confirmation. The registry labels the problem open; this is supplied status information, not an independently verified assessment.
\subsection{Sources}
\begin{itemize}
\item Supplied local UnsolvedMath registry record: permanent ID \(700923\), exact source code \texttt{catalogue entry}, title \emph{Kourovka Notebook Problem 13.5}; upstream revision \texttt{d371ee917ac9647fbb06cd6fc529deec6f3dda82}. The record was created and updated on 2026-10-16.
\item \emph{The Kourovka Notebook (Unsolved Problems in Group Theory)}, original problem 13.5, 13th issue (1995), proposed by V. V. Belyaev. The supplied background cites the 21st issue edition (2026), \texttt{arXiv:1401.0300v48}.
\item The supplied collection description attributes editing of Issues 1--21 (1965--2026) to E. I. Khukhro and V. D. Mazurov. These bibliographic details and the registry's open label are supplied information and have not been independently verified.
\end{itemize}
\end{document}
